\begin{abstract}
Bài báo trình bày LogiGraph, hệ thống định tuyến giao hàng và quy hoạch mạng lưới kho trên dữ liệu bản đồ thật của 24 phường
nội thành Hà Nội lấy từ OpenStreetMap (38\,604 đỉnh, 69\,411 cạnh có hướng). Toàn bộ cấu trúc dữ liệu lõi được tự cài đặt
chỉ bằng mảng, không dùng bảng băm hay hàng đợi dựng sẵn: đồ thị có trọng số biểu diễn bằng Adjacency List; hàng đợi ưu tiên
thích ứng có Locator phục vụ Dijkstra và tính lại tuyến khi có tắc hoặc cấm đường; Union-Find với union by rank và path
compression phục vụ Kruskal, kiểm tra liên thông và phân cụm giao hàng; quy hoạch động cho bài toán xếp hàng 0/1 Knapsack và
TSP bất đối xứng (Held-Karp). Mỗi giải thuật được kiểm tra chéo với một lời giải chuẩn bằng 81 test tự động. Thực nghiệm trên
đồ thị tới 50\,000 đỉnh cho thấy: Union-Find giữ khoảng 1{,}2 bước mỗi thao tác từ $10^3$ đến $10^6$ phần tử; Bellman-Ford
và Dijkstra dùng mảng chậm hơn Dijkstra dùng heap từ 13 tới hơn 146 lần; và lợi thế của hàng đợi thích ứng so với heap lazy
phụ thuộc mật độ đồ thị, rõ rệt trên đồ thị dày nhưng không xuất hiện trên mạng đường gần phẳng. Quy tắc bỏ qua sự cố nằm ngoài
tuyến giúp phản hồi trong 0{,}4 ms thay vì khoảng 110 ms khi phải tính lại.
\end{abstract}

\begin{IEEEkeywords}
Dijkstra, Adaptable Priority Queue, Union-Find, Kruskal, Knapsack, Held-Karp, định tuyến đô thị, OpenStreetMap
\end{IEEEkeywords}
